\documentclass[11pt,a4paper]{article}
\usepackage[utf8]{inputenc}
\usepackage[T1]{fontenc}
\usepackage[portuguese]{babel}
\usepackage{geometry}
\geometry{margin=2.5cm}
\usepackage{listings}
\usepackage{xcolor}
\usepackage{amsmath}
\usepackage{amssymb}
\usepackage{hyperref}
\usepackage{enumitem}
\usepackage{tcolorbox}
\usepackage{tabularx}
\usepackage{booktabs}

\definecolor{codebg}{rgb}{0.95,0.95,0.95}
\definecolor{codegreen}{rgb}{0.1,0.5,0.1}
\definecolor{codepurple}{rgb}{0.5,0.1,0.5}
\definecolor{codegray}{rgb}{0.5,0.5,0.5}

\lstdefinestyle{python}{
    backgroundcolor=\color{codebg},
    commentstyle=\color{codegray}\itshape,
    keywordstyle=\color{codepurple}\bfseries,
    stringstyle=\color{codegreen},
    basicstyle=\ttfamily\small,
    breaklines=true,
    frame=single,
    language=Python,
    showstringspaces=false,
    tabsize=4,
    literate={á}{{\'a}}1 {ã}{{\~a}}1 {é}{{\'e}}1 {í}{{\'i}}1 {ó}{{\'o}}1 {õ}{{\~o}}1 {ú}{{\'u}}1 {ç}{{\c{c}}}1 {Á}{{\'A}}1 {Ã}{{\~A}}1 {É}{{\'E}}1 {Í}{{\'I}}1 {Ó}{{\'O}}1 {Õ}{{\~O}}1 {Ú}{{\'U}}1 {Ç}{{\c{C}}}1 {â}{{\^a}}1 {ê}{{\^e}}1 {ô}{{\^o}}1 {û}{{\^u}}1 {Â}{{\^A}}1 {Ê}{{\^E}}1 {Ô}{{\^O}}1 {Û}{{\^U}}1 {à}{{\`a}}1 {è}{{\`e}}1 {ì}{{\`i}}1 {ò}{{\`o}}1 {ù}{{\`u}}1 {ä}{{\"a}}1 {ö}{{\"o}}1 {ü}{{\"u}}1 {Ä}{{\"A}}1 {Ö}{{\"O}}1 {Ü}{{\"U}}1
}
\lstset{style=python}

\newtcolorbox{solucao}{
  colback=yellow!5!white,
  colframe=orange!60!black,
  title=Solução proposta,
  fonttitle=\bfseries
}

\title{\textbf{Data Analysis with Python 2024--2025}\\Parte 5: Pandas (Parte 3)\\\large Soluções dos Exercícios}
\author{Soluções independentes elaboradas para fins de estudo, com base nos enunciados originais}
\date{}

\begin{document}
\maketitle

\section*{Nota Importante}
Este material é uma adaptação das soluções independentes para os exercícios baseados no curso \textit{Data Analysis with Python 2024-2025}, oferecido pela \textbf{The University of Helsinki MOOC Center}. As soluções abaixo não são o gabarito oficial do curso.

\tableofcontents
\newpage

\section{Soluções - Pandas (Parte 3)}

\subsection*{Exercício 5.1 -- Divisão de Data Continua}
\begin{solucao}
\begin{lstlisting}
import pandas as pd

def split_date_continues():
    df = pd.read_csv("src/Helsingin_pyorailijamaarat.csv", sep=";")
    df = df.dropna(axis=0, how="all")
    df = df.dropna(axis=1, how="all")
    
    # Criando o DataFrame das datas
    dates = df["Päivämäärä"].str.split(expand=True)
    dates.columns = ["Weekday", "Day", "Month", "Year", "Hour"]
    
    # Dicionario para normalizar o mes
    month_map = {"tammi": 1, "helmi": 2, "maalis": 3, "huhti": 4, "touko": 5, "kesä": 6, 
                 "heinä": 7, "elo": 8, "syys": 9, "loka": 10, "marras": 11, "joulu": 12}
    dates["Month"] = dates["Month"].map(month_map)
    dates["Year"] = dates["Year"].astype(int)
    dates["Day"] = dates["Day"].astype(int)
    dates["Hour"] = dates["Hour"].str.split(":", expand=True)[0].astype(int)
    
    df = df.drop(columns=["Päivämäärä"])
    return pd.concat([dates, df], axis=1)
\end{lstlisting}
Esta solução retoma o processamento de texto feito nas partes anteriores, e concatena o novo DataFrame de colunas temporais com as medições limpas usando \texttt{axis=1}.
\end{solucao}

\subsection*{Exercício 5.2 -- Clima de Ciclismo}
\begin{solucao}
\begin{lstlisting}
def cycling_weather():
    cycling = split_date_continues()
    weather = pd.read_csv("src/kumpula-weather-2017.csv")
    
    merged = pd.merge(cycling, weather, left_on=["Year", "Month", "Day"], right_on=["Year", "m", "d"])
    merged = merged.drop(columns=["m", "d", "Time", "Time zone"])
    return merged
\end{lstlisting}
O método \texttt{pd.merge} é acionado usando as chaves de ano, mês e dia. Como no DataFrame climático o mês e o dia são denotados por 'm' e 'd', informamos ao pandas essas referências diferentes utilizando \texttt{left\_on} e \texttt{right\_on}.
\end{solucao}

\subsection*{Exercício 5.3 -- Top Bandas}
\begin{solucao}
\begin{lstlisting}
def top_bands():
    top40 = pd.read_csv("src/UK-top40-1964-1-2.tsv", sep="\t")
    bands = pd.read_csv("src/bands.tsv", sep="\t")
    
    merged = pd.merge(top40, bands, left_on=["Artist"], right_on=["Band"])
    return merged
\end{lstlisting}
\end{solucao}

\subsection*{Exercício 5.4 -- Ciclistas por Dia}
\begin{solucao}
\begin{lstlisting}
import matplotlib.pyplot as plt

def cyclists_per_day():
    df = split_date_continues()
    df = df.drop(columns=["Hour", "Weekday"])
    return df.groupby(["Year", "Month", "Day"]).sum()

def main():
    df_daily = cyclists_per_day()
    
    # Restringir a Agosto (Mes=8) do ano 2017
    august = df_daily.loc[(2017, 8, slice(None))]
    
    # Opcional: ajustar o indice para ter apenas o Dia para o eixo x
    august = august.reset_index(level=[0, 1], drop=True)
    
    august.plot(figsize=(10,6))
    plt.xlabel("Dia")
    plt.xticks(range(1, 32))
    plt.show()

if __name__ == "__main__":
    main()
\end{lstlisting}
Através do \texttt{groupby} junto do \texttt{sum()}, obtivemos o somatório das observações diárias e o DataFrame ficou com um índice múltiplo (Ano, Mês, Dia). O \texttt{loc} permite acessar o recorte de 2017 e Mês 8 com \texttt{slice(None)} abrangendo todos os dias.
\end{solucao}

\subsection*{Exercício 5.5 -- Melhor Gravadora}
\begin{solucao}
\begin{lstlisting}
def best_record_company():
    df = pd.read_csv("src/UK-top40-1964-1-2.tsv", sep="\t")
    
    # Encontrar a gravadora que soma o maior valor na coluna WoC
    best_publisher = df.groupby("Publisher")["WoC"].sum().idxmax()
    
    return df[df["Publisher"] == best_publisher]
\end{lstlisting}
Com \texttt{df.groupby("Publisher")["WoC"].sum()} obtemos a soma de semanas para cada gravadora. A aplicação encadeada de \texttt{idxmax()} retorna a chave (nome da gravadora) com o maior valor.
\end{solucao}

\subsection*{Exercício 5.8 -- Séries Temporais de Bicicletas}
\begin{solucao}
\begin{lstlisting}
def bicycle_timeseries():
    df = split_date_continues()
    
    # Criar string unificada de data e hora
    datetime_str = df["Year"].astype(str) + "-" + df["Month"].astype(str) + "-" + \
                   df["Day"].astype(str) + " " + df["Hour"].astype(str) + ":00:00"
                   
    df["Date"] = pd.to_datetime(datetime_str)
    df = df.set_index("Date")
    df = df.drop(columns=["Year", "Month", "Day", "Hour", "Weekday"])
    
    return df
\end{lstlisting}
Aproveitamos as informações divididas no exercício 5.1 para juntá-las num formato reconhecível por \texttt{pd.to\_datetime}. Em seguida, usamos \texttt{set\_index("Date")} e varremos fora as colunas individuais temporais.
\end{solucao}

\subsection*{Exercício 5.9 -- Deslocamento (Commute)}
\begin{solucao}
\begin{lstlisting}
def commute():
    df = bicycle_timeseries()
    
    # Filtrar Agosto de 2017
    august = df["2017-08-01":"2017-08-31"]
    
    # Adicionar novamente a coluna de dias da semana (onde 1=Segunda, 7=Domingo)
    august["Weekday"] = august.index.dayofweek + 1
    
    # Agrupar pelo dia da semana e agregar pela soma
    commute_df = august.groupby("Weekday").sum()
    
    return commute_df

def main():
    commute_df = commute()
    commute_df.plot(figsize=(8,5))
    weekdays = ["Seg", "Ter", "Qua", "Qui", "Sex", "Sab", "Dom"]
    plt.xticks(range(1, 8), weekdays)
    plt.show()

if __name__ == "__main__":
    main()
\end{lstlisting}
No \texttt{DatetimeIndex}, \texttt{dayofweek} retorna valores de 0 a 6. Somando 1, formatamos o recorte entre 1 (Segunda) e 7 (Domingo). O \texttt{groupby} fará a união estatística somando as contagens diárias de todos os dias compatíveis de Agosto.
\end{solucao}

\vspace{2em}
\noindent\textbf{Créditos:} Este conteúdo foi originalmente criado pela \textit{The University of Helsinki MOOC Center} para o curso \textit{Data Analysis with Python}.
\end{document}
